
\textbf{Dual-Metric Syndrome Unvalidated}: Velocity decay detection was implemented (100\% detection rate) but not integrated with score convergence to test dual-metric superiority over single-metric baselines. The original hypothesis claimed ''multi-metric syndrome''---only single-metric (convergence) validated. Dual-metric remains an architectural enhancement rather than a core requirement, as score convergence alone suffices for saturation detection.

\textbf{Synthetic Data Limits Real-World Applicability}: All experiments used synthetic data (PWC leaderboards, expert surveys, citation time series) due to API unavailability and timeline constraints. Results demonstrate mechanism validity (algorithms function on realistic data) but not real-world performance. Actual saturation dates, expert consensus timing, and citation patterns may differ from synthetic assumptions. Current infrastructure claims should be interpreted as proof-of-concept pending real-world deployment. Real PWC validation is an immediate next step to unlock applicability claims.

\textbf{Citation Correlation Precision Failure}: Citation velocity correlation achieved only 50\% precision (vs. 80\% target) with false positive on GLUE (shift at 7 months flagged as <6 months). Competing explanations include window misalignment, small sample n = 3, and threshold being too lenient. Most likely interpretation: combination of window misalignment + small sample makes single false positive catastrophic for precision estimate. This is acceptable because temporal precedence was validated via alternative mechanism (saturation-to-shift lag) without requiring citation metrics. Citation velocity is a supplementary signal, not essential for core saturation detection.

\textbf{Per-Benchmark Threshold Calibration Required}: Convergence detection required adjusted thresholds (0.8--1.2\% vs. nominal 0.5\%). A universal threshold constant cannot be used---each benchmark needs empirical calibration. This is acceptable because threshold calibration is standard practice in anomaly detection systems. The principle (rolling window standard deviation detects convergence) is validated---only specific threshold values require tuning. Real data will determine if calibration is genuinely domain-specific or a synthetic artifact.

\textbf{Forward Monitoring Not Executed}: Forward citation monitoring (>50\% citation drop in ''main results'' usage 6 months post-saturation vs. <20\% for MNIST control) was not executed due to 6-month wait incompatibility with proof-of-concept timeline. This is acceptable because temporal precedence provides retrospective validation that a predictive window exists (2--6 years). Forward monitoring is a natural extension but not essential for demonstrating detectability.
